\section{Branch-tree records and the runtime qualification boundary}
\label{sec:b-data}

\subsection{Collection through the real fork path}

One root is generated as follows.  The prefix runs once under the frozen agent
to checkpoint 1 and is captured.  The live declared state is quiesced and
manifested (state surface, hashes, versions, budget counters), and isolated
sibling states are created immediately.  The verify-now tail and the continue
tail are then executed; on the continue branch, checkpoint 2 is captured when
reached and forks into verify-at-2 and continue.  Branch order and seed coupling
follow a predeclared rule.  Final artifacts are adjudicated by a hidden,
regime-blind evaluator on the registered horizon, and append-only records are
written.  The tree is executed through the actual adapter/fork path
(\cref{sec:b-experiments}); supplied leaf constants are not an acceptable
substitute for executed outcomes.

\subsection{The \texttt{b\_branch\_v2} record}
\label{sec:b-branch-schema}

Every node carries a schema version, paper/study/protocol hash, run/template/root
identifiers, split, checkpoint id and rule version, and an evaluator-only
condition label.  Every branch additionally carries:
\begin{itemize}[leftmargin=1.4em]
\item \texttt{decision\_index} and \texttt{parent\_branch\_id};
\item \texttt{history\_hash\_at\_decision}, so a checkpoint-2 state transplanted
  from another branch is detectable and rejected;
\item \texttt{remaining\_budget} at the decision;
\item checkpoint-2 features and feasibility mask, recorded only when reached;
\item \texttt{reached} and \texttt{terminated\_before\_decision};
\item repetition identifiers per decision;
\item the cost-ledger entry;
\item execution status, failure reason and \texttt{integrity\_status}.
\end{itemize}
All descendants inherit the root, template family and split of their origin.
The learner dataset is a restricted projection of the record; the outcome table
and the evaluator-only metadata are unreachable through feature construction,
and a leak test must fail when the projection is given access to either.  The
converter from records to tree recomputes the three regimes and rejects a
checkpoint-2 state transplanted from another branch, a denominator conditioned
on both policies reaching checkpoint 2, lineage or split leakage, and a second
verification in one root.

\subsection{The qualification gate (B-Q)}

A separate engineering pool with known labels is used.  Required checks:
identical serialized starting payloads where equality is claimed; immutable
read-only inputs and isolated writable state; no cross-branch mutable links,
queues or counters; pre-treatment mask and feature construction; stable
action/continuation/adjudicator versions; branch-order and seed-handling checks;
deterministic known-answer controls and stochastic-law diagnostics where
relevant; explicit missing-state and tool-failure handling; and fork
reachability and integrity.  Tolerances are matched to the state contract before
the scientific pilot.  A qualification check is evidence about declared
components, not proof of all hidden dependencies; unresolved outcome-relevant
state blocks causal wording for that substrate.  Only after B-Q passes for a
specific runtime/configuration may new branch contrasts be interpreted
scientifically.

\subsection{Splits}

Qualification, screen, pilot and final pools are disjoint at the template-family
level.  A single registry spans both the one-step and two-checkpoint lanes: a
template family used for training, validation or calibration cannot enter the
final test, and screen and pilot families are disjoint.  Training fits
representations and learners; validation selects architecture and candidate
policies; calibration fixes thresholds and the frozen controller; the final test
evaluates the frozen controller.  No test-time feature may contain future
outcomes, treatment effects or injector labels.
